% STATUS: written 2026-11-06 (English gate clean; sources read in full as recorded below).
% Sources read in full: Hochschild 1945 (Ann. of Math. 46, 58--67), Hochschild 1942 (Amer. J. Math. 64, 677--694).
\section[Derivations as one-cocycles: Hochschild 1945][Derivations as one-cocycles]{Derivations as one-cocycles: Hochschild 1945}\label{sec:origins:cohomology}

We now come to the end of the story. In 1945 Hochschild placed derivations inside a cohomology theory. His paper
is \emph{On the cohomology groups of an associative
algebra}~\cite{Hochschild1945}.

Hochschild states his aim at once. The paper is
``concerned primarily with the connections between the structure of an
algebra and the vanishing of its cohomology groups''~\cite[p.~58]{Hochschild1945}.
His model is a pair of results on Lie algebras, the lemmas of Whitehead. The
first lemma is about derivations of a semisimple Lie algebra into a module.
Hochschild looks for the same kind of result for associative algebras. He
finds the answer in the separable algebras. These are exactly the algebras
on which every derivation is inner. He also notes a second link. The
two-dimensional groups describe the extensions of $A$. These are the
algebras that have $A$ as a homomorphic image~\cite[p.~58]{Hochschild1945}.

His definitions are as follows. Let $A$ be an associative algebra over a field.
Let $M$ be an $A$-bimodule, that is, a two-sided module. An $m$-cochain is a
linear map from the $m$-fold tensor power of $A$ into
$M$~\cite[\S 1, p.~59]{Hochschild1945}. Hochschild defines a coboundary
operator $\delta$. It sends $m$-cochains to
$(m+1)$-cochains~\cite[\S 2, Definition~2.1, p.~59]{Hochschild1945}. He
credits the coboundary to an unpublished definition of Eilenberg and
Mac~Lane. He notes that dimensions one and two had appeared before. They
appear in the first and second lemmas of
Whitehead~\cite[p.~58]{Hochschild1945}.

In the lowest dimensions the definitions are short. A $0$-cochain is an
element $u \in M$, and
\[
(\delta u)(a) = a\cdot u - u\cdot a .
\]
For a $1$-cochain $f$,
\[
(\delta f)(a_1, a_2) = a_1 \cdot f(a_2) - f(a_1 a_2) + f(a_1)\cdot a_2
\]
\cite[p.~60]{Hochschild1945}. A cochain $f$ is a cocycle if $\delta f = 0$.
It is a coboundary if $f = \delta g$ for some cochain $g$. The
$m$-dimensional cohomology group $H^m(A,M)$ is the group of $m$-cocycles
modulo the coboundaries~\cite[p.~60, Definitions~2.2, 2.3]{Hochschild1945}.

In modern language, a $1$-cocycle is a linear map
$f\colon A \to M$ with $f(a_1a_2) = a_1 f(a_2) + f(a_1) a_2$. So a
$1$-cocycle is a derivation of $A$ into $M$. A $1$-coboundary is a map
$a \mapsto a u - u a$. So a $1$-coboundary is an inner derivation. Hence
$H^1(A,M)$ is the group of derivations modulo the inner ones.

With these definitions, Hochschild reaches his main results. He says that the theory is ``in a sense''
``degenerate''. For every $M$ there is another $A$-bimodule, written $(A:M)$.
For $m \ge 2$ it gives
$H^m(A,M) \cong H^{m-1}(A, (A:M))$~\cite[Theorem~3.1, p.~60]{Hochschild1945}.
So the groups of dimension one already determine all the others. In other
words, derivations control the whole theory.

\begin{theorem}[{\cite[p.~58 and \S 4]{Hochschild1945}}]\label{thm:origins:hochschild-separable}
An algebra is separable if and only if all its cohomology groups vanish.
Equivalently, it is separable if and only if every derivation of it into
every two-sided module is inner.
\end{theorem}

This theorem extends an earlier result of Hochschild himself. The earlier
result needed a non-modular ground field. A footnote places it in
Theorems~3.2 and~3.4 of his paper on semisimple algebras and generalized
derivations~\cite[p.~58, footnote~3]{Hochschild1945}.

That earlier paper begins with this sentence: ``Roughly speaking, a derivation of an algebra is
the infinitesimal operation corresponding to an
automorphism''~\cite[p.~677]{Hochschild1942}. The definition there is the one
of Jacobson from 1937. The first example is $\exp(tD)$, which acts on
polynomials as a one-parameter group of automorphisms.

So the lines of this chapter meet here. Leibniz gave the product rule.
Dedekind, Weber and Steinitz used it inside fields. Born, Jordan, Dirac, Ore
and Jacobson carried it to noncommutative algebras. On the side of groups, the
same objects are the crossed homomorphisms of
Section~\ref{sec:origins:crossed-homomorphisms}. In the work of Hochschild
all of these become one notion. The rest of the book starts from here.
